
This work presents proof-of-concept validation of executable API contracts for ML reproducibility. We implemented and evaluated three contract types: structural contracts (import-time validation of shapes and dtypes), metamorphic contracts (runtime probes for mathematical properties), and composition contracts (cross-library consistency validation).

Our controlled experiments demonstrate technical feasibility:
\begin{itemize}
\item Structural contracts detected 2/2 real scenarios at import time with <0.03s overhead
\item Metamorphic contracts detected 40/40 synthetic violations with 0/10 false positives in 3.7ms
\item Composition contracts detected 24/24 synthetic defects with 0/6 false positives in <1ms
\end{itemize}

These results validate the hypothesis that certain classes of environment-stage defects are detectable via lightweight behavioral contracts. However, the small sample sizes (N=2, N=50, N=30), synthetic defect injection, and absence of version stability testing limit generalizability. The gap between proof-of-concept validation and production deployment is substantial.

\textbf{Critical Next Steps}: Validate contracts on real defect corpora (Jiang et al.'s 348 bugs or equivalent), test version stability across library updates, scale experiments to $N\geq 100$ per contract type, and conduct pilot deployments in production repositories. Without this additional validation, claims about practical applicability to ML reproducibility workflows remain speculative.

The path from proof-of-concept to production-ready tooling requires:
\begin{enumerate}
\item Statistical validation on real defects (not synthetic)
\item Version stability assessment across library updates
\item False positive rate measurement in production repositories
\item Adoption friction analysis via user studies
\end{enumerate}

If these validations succeed, API contracts could provide a complementary reproducibility practice alongside environment isolation, dependency pinning, and integration testing. The proof-of-concept results establish technical feasibility; extensive empirical validation is required to establish practical utility.
